% Generated draft: requires mathematical and source review.
\documentclass[11pt]{article}
\usepackage[margin=25mm]{geometry}
\usepackage{fontspec}
\setmainfont{Latin Modern Roman}
\usepackage{amsmath,amssymb,mathtools,unicode-math}
\setmathfont{Latin Modern Math}
\usepackage{enumitem,needspace,xurl,hyperref}
\setlength{\parindent}{0pt}
\setlength{\parskip}{5pt}
\setlength{\emergencystretch}{3em}
\newcommand{\N}{\mathbb{N}}
\newcommand{\Z}{\mathbb{Z}}
\newcommand{\Q}{\mathbb{Q}}
\newcommand{\R}{\mathbb{R}}
\newcommand{\C}{\mathbb{C}}
\newcommand{\F}{\mathbb{F}}
\newcommand{\A}{\mathbb{A}}
\newcommand{\PP}{\mathbb{P}}
\newcommand{\E}{\mathbb{E}}
\newcommand{\eps}{\varepsilon}
\newcommand{\dd}{\,\mathrm{d}}
\begin{document}
\section*{UnsolvedMath ID: 801040}
\textbf{Source problem code:} catalogue entry
\subsection{Definitions and mathematical statement}
Let \(X\) be a standard Borel space. A Borel equivalence relation on \(X\) is an equivalence relation \(E\subseteq X\times X\) that is Borel. It is countable if every equivalence class is countable. It is \emph{hyperfinite} if there exists a sequence \((E_n)_{n\in\N}\) of Borel equivalence relations on \(X\), each having only finite equivalence classes, such that
\[
E_n\subseteq E_{n+1}\quad(n\in\N),\qquad E=\bigcup_{n\in\N}E_n.
\]
Here the finite classes need not have a common finite bound on their sizes.

\textbf{Source question.} Suppose that there is a counterexample to the problem called the \emph{hf/hf problem} in the source. Must the set of all hyperfinite equivalence relations be \(\boldsymbol{\Sigma}^{1}_{2}\)-complete?

\textbf{Unresolved formulation data.} The supplied record does not define the hf/hf problem, so its counterexample hypothesis cannot be expanded into explicit mathematical assumptions. It also does not specify how equivalence relations are coded, the ambient space of codes, or the reductions intended by completeness. Thus the displayed source question does not yet determine a fully specified descriptive-set-theoretic statement.

More explicitly, a coded formulation would require a Polish parameter space \(P\), a specified subset \(C\subseteq P\) of admissible codes, and an interpretation assigning to every \(c\in C\) a Borel equivalence relation \(E_c\) on a standard Borel space \(X_c\). One could then define
\[
\mathcal H=\{c\in C:E_c\text{ is hyperfinite}\}\subseteq P.
\]
For a Polish space \(Y\), a set \(A\subseteq Y\) belongs to \(\boldsymbol{\Sigma}^{1}_{2}(Y)\) if there is a Borel set \(B\subseteq Y\times\N^{\N}\times\N^{\N}\) such that
\[
y\in A\iff (\exists u\in\N^{\N})(\forall v\in\N^{\N})\ (y,u,v)\in B.
\]
After specifying a class \(\mathscr R\) of reduction maps from \(\N^{\N}\) to \(P\), the intended completeness assertion would require
\[
\mathcal H\in\boldsymbol{\Sigma}^{1}_{2}(P)
\quad\text{and}\quad
(\forall A\in\boldsymbol{\Sigma}^{1}_{2}(\N^{\N}))
(\exists r\in\mathscr R)(\forall x\in\N^{\N})\,
\bigl(x\in A\iff r(x)\in\mathcal H\bigr).
\]
This is a formulation schema only: neither the coding data nor the reduction class is supplied. In particular, continuous and Borel reductions are not silently identified. The definition of hyperfiniteness above supplies standard terminology, but the source's intended domain of equivalence relations still requires confirmation.
\subsection{Short English statement}
The source asks whether the existence of a counterexample to the hf/hf problem forces the collection of hyperfinite equivalence relations to be \(\boldsymbol{\Sigma}^{1}_{2}\)-complete. A precise formulation requires the definition of hf/hf, the intended coding of equivalence relations, and the reduction convention; these are absent from the supplied record.
\subsection{Sources}
\begin{itemize}
\item Supplied local UnsolvedMath registry record: permanent ID \(801040\), exact source problem code \texttt{catalogue entry}, title \emph{Open Problems in Mathematical Logic --- Problem 52}. The registry attributes the problem to Barmpalias, Gao et al., with proposed year 2026, and labels it open. This is supplied registry information, not an independently verified status assessment.
\item Supplied source list: Barmpalias, Gao et al., \emph{Open Problems in Mathematical Logic} (2026), Problem 52, \url{https://arxiv.org/abs/2608.26628}. The record reports access on 2026-10-16 and extraction by \texttt{arxiv\_latex}. The linked document was not independently consulted for this draft.
\item The statement describes the question as inspired by the Frisch--Shinko--Vidny\'anszky theorem. No theorem statement or separate bibliographic reference for that theorem is supplied.
\item Supplied upstream identity: \(801040\); revision \texttt{d371ee917ac9647fbb06cd6fc529deec6f3dda82}. The record gives creation and update timestamps of \texttt{2026-10-16T14:00:00Z}.
\end{itemize}
\end{document}
